% =====================================================================
% di_style.tex: Differential Equations homework / test layout.
% Loads the house preamble (local copies of the AoI worksheet files:
% ws_preamble_build.tex, ws_preamble_local.tex -> ws_figstyles.tex),
% then adds the homework layout used by HW1-HW3 and the Test.
% Usage (before \begin{document}):  \ditle{Homework 1}
% =====================================================================
% WORKSHEET COPY of ../notes/preamble_build.tex (2026-10-09 version), unchanged.
% preamble_build.tex: copy of Maximus's preamble.tex for the DE v3 build.
% The blocks marked [BUILD COPY CHANGE] are the earlier build fixes.
% The blocks marked [V3 ADDITION] are new in this version (example box,
% spacing). Everything else is unchanged from preamble.tex.

% For margins
\usepackage[a4paper, margin = 20mm]{geometry}

\usepackage[Lenny]{fncychap}
% Rejne, Glenn, Lenny

\usepackage{etoc}

% [V3 ADDITION] cancel: used to show the dx's "cancelling" in Section 11.3
\usepackage{cancel}
\renewcommand{\CancelColor}{\color{red}}

% For enumeration & math symbols
\usepackage{enumerate}
\usepackage{amsmath,amsfonts,amssymb,amsthm,graphicx,lipsum}

% For diagrams, plots and boxes
\usepackage{pgfplots}
\pgfplotsset{compat=1.18}
\usepackage{tikz}
\usepackage[most]{tcolorbox}
\tcbuselibrary{skins}
\usepackage{xcolor}
\pgfplotsset{
    standard/.style={clip = false,
    axis line style = thick,
    trig format=rad,
    axis lines = middle,
    every axis x label/.style={at={(current axis.right of origin)},anchor= west},
    every axis y label/.style={at={(current axis.above origin)},anchor=south}
    }
}

% For GC font
\usepackage{fontspec}
% [BUILD COPY CHANGE 1] The TI-84 key font is loaded only if the file exists
% (checked in ./Fonts/... and ../Fonts/...). Otherwise \ti{} falls back to a
% small boxed sans-serif key. Known letters of the TI font are mapped to their
% key names in the fallback: o = [y=], z = [mode], f = [L3].
\newcommand{\tikey}[1]{\tcbox[on line, enhanced, boxsep=0pt, left=2pt, right=2pt,
  top=1pt, bottom=1pt, boxrule=0.5pt, arc=1.5pt, colback=white,
  colframe=black!70]{\sffamily\footnotesize\strut #1}}
\ExplSyntaxOn
\NewDocumentCommand{\tifallback}{m}{
  \str_case:nnF {#1}
    { {o}{\tikey{y=}} {z}{\tikey{mode}} {f}{\tikey{L3}} }
    { \tikey{#1} }
}
\ExplSyntaxOff
\def\tifontcase{0}
\IfFileExists{Fonts/TI-84 Plus CE Key Fonts/TI84PlusCEKeys.ttf}
  {\def\tifontcase{1}}{}
\IfFileExists{../Fonts/TI-84 Plus CE Key Fonts/TI84PlusCEKeys.ttf}
  {\def\tifontcase{2}}{}
\ifcase\tifontcase\relax
  % 0: font not found, use boxed fallback keys
  \newcommand{\ti}[1]{\tifallback{#1}}
\or
  % 1: original location (as in preamble.tex)
  \newfontfamily\tifont[
    Path = {Fonts/TI-84 Plus CE Key Fonts/},
    UprightFont = {TI84PlusCEKeys.ttf}
  ]{}
  \newcommand{\ti}[1]{{\tifont #1}}
\or
  % 2: one level up (when compiling inside draft-ch1/)
  \newfontfamily\tifont[
    Path = {../Fonts/TI-84 Plus CE Key Fonts/},
    UprightFont = {TI84PlusCEKeys.ttf}
  ]{}
  \newcommand{\ti}[1]{{\tifont #1}}
\fi

% [AoI items 17/18/37] GC DISPLAY FONT. Every GC key, command, menu item and
% screen line goes through \gcfont, so the font can be swapped HERE in one place.
%   \gc{fnInt(}      command / menu item / stored variable in the display font
%   \gckey{2nd}      a calculator key (display font in a small key outline)
%   \gcscreen{...}   a calculator screen (display font on a grey-green LCD panel)
% The font file lives in ../Fonts (shared by notes, slides and worksheet).
\def\gcfontfile{VT323-Regular.ttf}
\IfFileExists{../Fonts/\gcfontfile}{%
  \newfontfamily\gcfont[Path=../Fonts/, Scale=1.32]{\gcfontfile}}{%
  \newcommand\gcfont{\ttfamily}}
\newcommand{\gc}[1]{\mbox{\gcfont #1}}
\newcommand{\gckey}[1]{\tcbox[on line, enhanced, boxsep=0pt, left=2pt, right=2pt,
  top=0.6pt, bottom=0.2pt, boxrule=0.5pt, arc=1.5pt, colback=black!4,
  colframe=black!65]{\gcfont\strut #1}}
\definecolor{lcd}{RGB}{214,222,204}
\newcommand{\gcscreen}[1]{\tcbox[on line, enhanced, boxsep=1pt, left=4pt, right=4pt,
  top=2pt, bottom=2pt, boxrule=1.2pt, arc=1pt, colback=lcd, colframe=black!70]{%
  \gcfont\begin{tabular}{@{}l@{}}#1\end{tabular}}}
% integral template as shown on the TI-84 Plus CE (MathPrint): \gcint{0}{2}{e^(-X²)}{X}
\newcommand{\gcint}[4]{{\gcfont\raisebox{-0.35ex}{\Large ∫}\hspace{-0.15em}%
  \raisebox{-0.2ex}{\begin{tabular}{@{}l@{}}\footnotesize #2\\[-1.3ex]\footnotesize #1\end{tabular}}%
  \hspace{0.1em}(#3)\mathrm{d}#4}}
% old key macro kept as an alias so nothing breaks
\renewcommand{\ti}[1]{\gckey{#1}}

% SI units
\usepackage{siunitx}

% setspace
%\usepackage{setspace}
%\onehalfspacing

% For hyperlinks within and out of document
\usepackage{hyperref}
\hypersetup{
    colorlinks = true,
    linkcolor = black,
    urlcolor = blue,
}   

% For captions for pics
\usepackage{capt-of}


% tcolorbox custom environments
% Definition ========================================================
\newtcolorbox[auto counter,number within=section]{definition}[2][drop shadow]{%
enhanced, before skip=5mm, after skip=6mm, % [V3 ADDITION] spacing
colback=blue!5!white,colframe=blue!50!black,fonttitle=\bfseries,
title= Definition~\thetcbcounter. #2,#1}

% Theorem ==========================================================
\newtcolorbox[auto counter,number within=section]{theorem}[2][drop shadow]{%
enhanced, before skip=5mm, after skip=6mm, % [V3 ADDITION] spacing
colback=black!5!white,colframe=black,fonttitle=\bfseries,
title= Theorem~\thetcbcounter. #2,#1}

% Caution box ======================================================
\newtcolorbox{caution}[1][]{enhanced,
before skip=5mm,after skip=6mm, % [V3 ADDITION] was 2mm/3mm
boxrule=1pt,left=11mm,right=2mm,top=3mm,bottom=3mm,
colback=yellow!50,
colframe=yellow!20!black,
sharp corners,rounded corners=southeast,arc is angular,arc=3mm,
underlay={%
\path[fill=tcbcolback!80!black] ([yshift=3mm]interior.south east)--++(-0.4,-0.1)--++(0.1,-0.2);
\path[draw=tcbcolframe,shorten <=-0.05mm,shorten >=-0.05mm] ([yshift=3mm]interior.south east)--++(-0.4,-0.1)--++(0.1,-0.2);
\path[fill=yellow!50!black,draw=none] (interior.south west) rectangle node[white]{\Huge\bfseries !} ([xshift=9mm]interior.north west);
},
drop fuzzy shadow,#1}

% Question box ======================================================
\newtcolorbox{question}[1][]{enhanced,
before skip=5mm,after skip=6mm, % [V3 ADDITION] was 2mm/3mm
boxrule=1pt,left=11mm,right=2mm,top=3mm,bottom=3mm,
colback=blue!20!white,
colframe=blue!20!black,
sharp corners,rounded corners=southeast,arc is angular,arc=3mm,
underlay={%
\path[fill=tcbcolback!80!black] ([yshift=3mm]interior.south east)--++(-0.4,-0.1)--++(0.1,-0.2);
\path[draw=tcbcolframe,shorten <=-0.05mm,shorten >=-0.05mm] ([yshift=3mm]interior.south east)--++(-0.4,-0.1)--++(0.1,-0.2);
\path[fill=blue!50!black,draw=none] (interior.south west) rectangle node[white]{\Huge\bfseries ?} ([xshift=9mm]interior.north west);
},
drop fuzzy shadow,#1}

% Information box ======================================================
\newtcolorbox{info}[1][]{enhanced,
before skip=5mm,after skip=6mm, % [V3 ADDITION] was 2mm/3mm
boxrule=1pt,left=11mm,right=2mm,top=3mm,bottom=3mm,
colback=cyan!20!white,
colframe=cyan!20!black,
sharp corners,rounded corners=southeast,arc is angular,arc=3mm,
underlay={%
\path[fill=tcbcolback!80!black] ([yshift=3mm]interior.south east)--++(-0.4,-0.1)--++(0.1,-0.2);
\path[draw=tcbcolframe,shorten <=-0.05mm,shorten >=-0.05mm] ([yshift=3mm]interior.south east)--++(-0.4,-0.1)--++(0.1,-0.2);
\path[fill=cyan!50!black,draw=none] (interior.south west) rectangle node[white]{\Huge\bfseries i} ([xshift=9mm]interior.north west);
},
drop fuzzy shadow,#1}

% Mistake box ======================================================
\newtcolorbox{mistake}[1][]{enhanced,
before skip=5mm,after skip=6mm, % [V3 ADDITION] was 2mm/3mm
boxrule=1pt,left=11mm,right=2mm,top=3mm,bottom=3mm,
colback=red!20!white,
colframe=red!20!black,
sharp corners,rounded corners=southeast,arc is angular,arc=3mm,
underlay={%
\path[fill=tcbcolback!80!black] ([yshift=3mm]interior.south east)--++(-0.4,-0.1)--++(0.1,-0.2);
\path[draw=tcbcolframe,shorten <=-0.05mm,shorten >=-0.05mm] ([yshift=3mm]interior.south east)--++(-0.4,-0.1)--++(0.1,-0.2);
\path[fill=red!50!black,draw=none] (interior.south west) rectangle node[white]{\Huge\bfseries $\times$} ([xshift=9mm]interior.north west);
},
drop fuzzy shadow,#1}



% [V3 ADDITION] Worked example box ==================================
% Own counter, numbered within the section (Example 11.2.1, 11.2.2, ...).
% For per-subsection numbering (11.2.1.1) change `section' to `subsection'. Blue `question' boxes are kept for thought-provoking questions
% only; worked examples always use this box. Usage:
%   \begin{example} ...question text... \end{example}
%   \solution{...}
\newtcolorbox[auto counter,number within=section]{example}[1][]{%
enhanced, breakable=false, before skip=5mm, after skip=4mm,
colback=teal!5!white, colframe=teal!55!black, boxrule=0.8pt,
left=3mm, right=3mm, top=2mm, bottom=2mm,
sharp corners, rounded corners=northwest, arc=3mm,
fonttitle=\bfseries, coltitle=white,
title=Example~\thetcbcounter, attach boxed title to top left={yshift=-2mm, xshift=3mm},
boxed title style={colback=teal!55!black, sharp corners, rounded corners=northwest, arc=2mm},
drop fuzzy shadow, #1}

% [V3 ADDITION] Spacing between sections, subsections, parts, paragraphs =
% (fncychap's chapter heading is untouched; only \section and below.)
\makeatletter
\renewcommand\section{\@startsection{section}{1}{\z@}%
  {-2ex \@plus -.5ex \@minus -.2ex}{3.2ex \@plus .3ex}%
  {\normalfont\LARGE\bfseries\color{blue!50!black}}}
\renewcommand\subsection{\@startsection{subsection}{2}{\z@}%
  {-1.5ex \@plus -.5ex \@minus -.2ex}{2.6ex \@plus .3ex}%
  {\normalfont\Large\bfseries}}
\renewcommand\subsubsection{\@startsection{subsubsection}{3}{\z@}%
  {-4ex \@plus -1ex \@minus -.5ex}{1.6ex \@plus .2ex}%
  {\normalfont\large\bfseries\color{teal!50!black}}}
\makeatother
\setlength{\parindent}{0pt}
\setlength{\parskip}{0.9ex plus 0.3ex minus 0.2ex}

% Headers and Footers
\usepackage{fancyhdr,lastpage}
% [V3 R2] grey left-margin line on the out-of-syllabus pages (item 13)
\usepackage{eso-pic}
% [AoI item 63] thicker line (2.6pt, was 1.6pt) that starts just BELOW the header rule
% (top 14mm of the page left clear). Flag: copy this line into every chapter's preamble.
\newcommand{\oosmargin}{\AddToShipoutPictureBG{\put(12mm,0){\color{gray!75}\rule{2.6pt}{\dimexpr\paperheight-14mm\relax}}}}
\pagestyle{fancy}
\lhead{Maximus Pung Jun}
% [BUILD COPY CHANGE 2] Header bug fix: \chaptername is the word "Chapter",
% so the original header printed "Chapter 1: Chapter". \chaptermark now stores
% the bare chapter title in \leftmark, and the header uses \leftmark.
\renewcommand{\chaptermark}[1]{\markboth{#1}{}}
\rhead{Chapter \thechapter: \leftmark}
% \cfoot{Page \thepage\  of \pageref{LastPage}} % buggy, does not appear on chapter pages.

% [V3 R2] tick and cross markers for mistake boxes (G3)
\newcommand{\wrong}{\textcolor{red!70!black}{\ensuremath{\boldsymbol{\times}}}}
\newcommand{\cmark}{\textcolor{green!50!black}{\ensuremath{\checkmark}}}

\newcommand{\xmark}{\wrong}   % alias: red cross
% [V3 R2] out-of-syllabus helpers (DE notes item 13)
%   \oosmargin    : grey line in the left margin on every following page
%   \oosmarginoff : stop the margin line
%   \oosnote      : grey banner "Out of syllabus. For your interest only."
%   outofsyllabus : environment = new page + margin line + banner; ends with
%                   \clearpage and switches the margin line off again
\newcommand{\oosmarginoff}{\ClearShipoutPictureBG}
\newcommand{\oosnote}[1][Nothing in this section is examined in H2 Mathematics.]{\par\noindent\colorbox{gray!15}{\parbox{\dimexpr\textwidth-2\fboxsep}{\textbf{Out of syllabus. For your interest only.} #1}}\par}
\newenvironment{outofsyllabus}{\clearpage\oosmargin}{\clearpage\oosmarginoff}

% Notation as in Maximus's notes: italic d, display style everywhere (G1).
% Fractions inside exponents stay small, so never use these macros inside ^{}.
\newcommand{\dydx}{\dfrac{\mathrm{d}y}{\mathrm{d}x}}
\newcommand{\dvdx}{\dfrac{\mathrm{d}v}{\mathrm{d}x}}
\newcommand{\dudx}{\dfrac{\mathrm{d}u}{\mathrm{d}x}}
\newcommand{\ddydx}{\dfrac{\mathrm{d}^{2}y}{\mathrm{d}x^{2}}}
\newcommand{\ddt}[1]{\dfrac{\mathrm{d}#1}{\mathrm{d}t}}
% Small syllabus flag (no longer used in the DE notes, kept for other chapters)
\newcommand{\syllabusflag}[1]{\par\noindent{\footnotesize\color{red!70!black}\textbf{Syllabus note.} #1}\par}


% Teacher's copy toggle ================================================
% [BUILD COPY CHANGE 3] \ifx compares the whole word (the original
% \if\conform test was always false). DE_notes_student.tex predefines
% \conform as HideSolutions.
% A student copy defines \conform as HideSolutions BEFORE loading this preamble.
\providecommand\conform{ShowSolutions}
\newcommand\ShowSolutionsKey{ShowSolutions}
% [V3 ADDITION] In the student copy the hidden solution leaves stretchable
% blank space (\vfill), so the free space on each page is shared out as
% working space below the examples instead of collecting at the bottom.
% [AoI item 10] "Solution." on its own line, so that (a), (b), (c) line up below it.
\newcommand\solution[1]{\ifx\conform\ShowSolutionsKey
  \par\noindent\textbf{Solution.}\par\nobreak\vspace{0.5ex}\noindent #1\par\medskip\else\par\vspace{\workspace}\fi}
% [V3 R2] minimum working space left under each example in the student copy
\newlength\workspace\setlength\workspace{3cm plus 1fill}


\newcommand{\NN}{\mathbb{N}}
\newcommand{\ZZ}{\mathbb{Z}}
\newcommand{\QQ}{\mathbb{Q}}
\newcommand{\RR}{\mathbb{R}}
\newcommand{\CC}{\mathbb{C}}


% WORKSHEET COPY of ../notes/preamble_local.tex (version the worksheet was last built with,
% from ../_backup_pre_limitsum, 2026-10-09 17:31). Local so the worksheet no longer depends
% on ../notes; changes: % WORKSHEET COPY of ../notes/figstyles.tex (2026-10-09 17:31 version). Only change (follow-up):
% thin light axes in `sketch' and `ax3' (marked [2026-10-09 follow-up]).
% =====================================================================
% figstyles.tex: plot styles and figure macros shared by the notes, the
% slides (../slides) and the worksheets (../worksheet). Needs pgfplots,
% the fillbetween library and tikz libraries arrows.meta and calc.
% House rules built in:
%   * `axis on top' in every sketch: shading never covers an axis, and the
%     fill starts at the axis line (items 31, 32, 34);
%   * regions are filled from the curve itself (fill between / closed plot);
%   * 3D solids in the page-13 style (translucent blue, light end caps).
% =====================================================================
\pgfplotsset{
  sketch/.style={clip=false, axis line style={thin, black!45}, every axis label/.append style={black!60}, trig format=rad, % [2026-10-09 follow-up] was axis line style=thick
    axis lines=middle, axis on top,
    every axis x label/.style={at={(current axis.right of origin)},anchor=west},
    every axis y label/.style={at={(current axis.above origin)},anchor=south},
    ticks=none, xlabel=$x$, ylabel=$y$, width=7.6cm, height=6cm,
    every node/.append style={font=\footnotesize}},
  curve/.style={blue, semithick, smooth, samples=150},
  line2/.style={red!75!black, semithick, samples=2},
  % 3D solids of revolution (angles in DEGREES: no `trig format=rad' here)
  solid3d/.style={hide axis, trig format=deg, width=8.4cm, height=6.4cm, clip=false,
    z buffer=sort, colormap={bl}{rgb255=(150,185,235) rgb255=(40,90,180)},
    every node/.append style={font=\footnotesize}},
  solidsurf/.style={surf, shader=interp, opacity=0.75, point meta=y},
  red3d/.style={surf, shader=flat, draw=none, point meta=1,
    colormap={r}{rgb255=(215,50,50) rgb255=(215,50,50)}},
  orange3d/.style={surf, shader=interp, point meta=z,
    colormap={o}{rgb255=(250,210,160) rgb255=(220,120,30)}},
  cap3d/.style={surf, shader=flat, draw=none, point meta=1,
    colormap={c}{rgb255=(185,205,238) rgb255=(185,205,238)}},
  mini3d/.style={solid3d, width=2.6cm, height=2.4cm},
}
\tikzset{
  region/.style={fill=blue!22, draw=none},
  regionB/.style={fill=orange!35, draw=none},
  strip/.style={fill=red!55, draw=red!70!black, thin},
  asym/.style={black, dashed, thin},
  pt/.style={circle, fill=black, inner sep=1.1pt},
  dimarrow/.style={<->, >=Stealth, thin},
  ax3/.style={-{Stealth[length=1.5mm]}, thin, black!50}, % [2026-10-09 follow-up] was -{Stealth[length=2mm]}, thick
}
\definecolor{hl}{RGB}{30,110,200}

% ---------------------------------------------------------------------
% Mini 3D solids for SOLUTIONS (items 54, 62, 66, 74): solid only, no
% washer, every picture 1.5 cm tall; a "big minus small" pair is 4.4 cm wide.
%  \solX[x unit cm]{R(x)}{a}{b}{R(a)}{R(b)}{x-axis end}{y-axis top}{curve colour}
%     solid swept out by y = R(x), a <= x <= b, about the x-axis
%  \solY[y unit cm]{R(y)}{c}{d}{R(c)}{R(d)}{x-axis end}{y-axis top}{curve colour}
%     solid swept out by x = R(y), c <= y <= d, about the y-axis
%     (R is written in terms of the letter x, which here stands for y)
% ---------------------------------------------------------------------
\newcommand{\solX}[9][0.5]{%
\begin{tikzpicture}[baseline=(current bounding box.center)]
\begin{axis}[mini3d, x={(#1cm,-0.06cm)}, y={(0.12cm,0.09cm)}, z={(0cm,0.4cm)}]
  \addplot3[cap3d, domain=0:1, y domain=0:360, samples=2, samples y=40] ({#3}, {x*(#5)*cos(y)}, {x*(#5)*sin(y)});
  \addplot3[solidsurf, opacity=0.5, point meta=z, domain=#3:#4, y domain=0:360, samples=24, samples y=40] ({x}, {(#2)*cos(y)}, {(#2)*sin(y)});
  \addplot3[cap3d, domain=0:1, y domain=0:360, samples=2, samples y=40] ({#4}, {x*(#6)*cos(y)}, {x*(#6)*sin(y)});
  \addplot3[#9, thick, domain=#3:#4, samples=40, samples y=1] ({x}, {0}, {#2});
  \draw[ax3] (axis cs:{min(0,#3)-0.3},0,0) -- (axis cs:#7,0,0);
  \draw[ax3] (axis cs:0,0,{-0.15*(#8)}) -- (axis cs:0,0,#8);
\end{axis}
\end{tikzpicture}}
\newcommand{\solY}[9][0.6]{%
\begin{tikzpicture}[baseline=(current bounding box.center)]
\begin{axis}[mini3d, x={(0.3cm,-0.03cm)}, y={(0.13cm,0.11cm)}, z={(0cm,#1cm)}]
  \addplot3[cap3d, domain=0:1, y domain=0:360, samples=2, samples y=40] ({x*(#5)*cos(y)}, {x*(#5)*sin(y)}, {#3});
  \addplot3[solidsurf, opacity=0.5, point meta=z, domain=#3:#4, y domain=0:360, samples=24, samples y=40] ({(#2)*cos(y)}, {(#2)*sin(y)}, {x});
  \addplot3[cap3d, domain=0:1, y domain=0:360, samples=2, samples y=40] ({x*(#6)*cos(y)}, {x*(#6)*sin(y)}, {#4});
  \addplot3[#9, thick, domain=#3:#4, samples=40, samples y=1] ({#2}, {0}, {x});
  \draw[ax3] (axis cs:0,0,{min(0,#3)-0.15*(#8)}) -- (axis cs:0,0,#8);
  \draw[ax3] (axis cs:-0.3,0,0) -- (axis cs:#7,0,0);
\end{axis}
\end{tikzpicture}}
\newcommand{\minisolo}[1]{\adjustbox{totalheight=1.5cm, max width=4.6cm}{#1}}
\newcommand{\minipair}[2]{\adjustbox{totalheight=1.5cm, max width=4.6cm}{#1\;\raisebox{0.15ex}{\Large$-$}\;#2}}
\newcommand{\minilabel}[1]{\raisebox{0.65cm}{\small #1}\;}
; [2026-10-09 follow-up] \mk marks now bold, as in the notes.
% =====================================================================
% preamble_local.tex: extra packages, styles and macros for the
% Applications of Integration chapter ONLY. Loaded after preamble_build.tex.
% Plot styles and the mini 3D macros are in figstyles.tex (shared with the
% slides and worksheets).
% =====================================================================
\usepgfplotslibrary{fillbetween}
\usepackage{adjustbox} % mini 3D pictures (figstyles.tex)
\usetikzlibrary{arrows.meta, decorations.pathreplacing, calc}
% WORKSHEET COPY of ../notes/figstyles.tex (2026-10-09 17:31 version). Only change (follow-up):
% thin light axes in `sketch' and `ax3' (marked [2026-10-09 follow-up]).
% =====================================================================
% figstyles.tex: plot styles and figure macros shared by the notes, the
% slides (../slides) and the worksheets (../worksheet). Needs pgfplots,
% the fillbetween library and tikz libraries arrows.meta and calc.
% House rules built in:
%   * `axis on top' in every sketch: shading never covers an axis, and the
%     fill starts at the axis line (items 31, 32, 34);
%   * regions are filled from the curve itself (fill between / closed plot);
%   * 3D solids in the page-13 style (translucent blue, light end caps).
% =====================================================================
\pgfplotsset{
  sketch/.style={clip=false, axis line style={thin, black!45}, every axis label/.append style={black!60}, trig format=rad, % [2026-10-09 follow-up] was axis line style=thick
    axis lines=middle, axis on top,
    every axis x label/.style={at={(current axis.right of origin)},anchor=west},
    every axis y label/.style={at={(current axis.above origin)},anchor=south},
    ticks=none, xlabel=$x$, ylabel=$y$, width=7.6cm, height=6cm,
    every node/.append style={font=\footnotesize}},
  curve/.style={blue, semithick, smooth, samples=150},
  line2/.style={red!75!black, semithick, samples=2},
  % 3D solids of revolution (angles in DEGREES: no `trig format=rad' here)
  solid3d/.style={hide axis, trig format=deg, width=8.4cm, height=6.4cm, clip=false,
    z buffer=sort, colormap={bl}{rgb255=(150,185,235) rgb255=(40,90,180)},
    every node/.append style={font=\footnotesize}},
  solidsurf/.style={surf, shader=interp, opacity=0.75, point meta=y},
  red3d/.style={surf, shader=flat, draw=none, point meta=1,
    colormap={r}{rgb255=(215,50,50) rgb255=(215,50,50)}},
  orange3d/.style={surf, shader=interp, point meta=z,
    colormap={o}{rgb255=(250,210,160) rgb255=(220,120,30)}},
  cap3d/.style={surf, shader=flat, draw=none, point meta=1,
    colormap={c}{rgb255=(185,205,238) rgb255=(185,205,238)}},
  mini3d/.style={solid3d, width=2.6cm, height=2.4cm},
}
\tikzset{
  region/.style={fill=blue!22, draw=none},
  regionB/.style={fill=orange!35, draw=none},
  strip/.style={fill=red!55, draw=red!70!black, thin},
  asym/.style={black, dashed, thin},
  pt/.style={circle, fill=black, inner sep=1.1pt},
  dimarrow/.style={<->, >=Stealth, thin},
  ax3/.style={-{Stealth[length=1.5mm]}, thin, black!50}, % [2026-10-09 follow-up] was -{Stealth[length=2mm]}, thick
}
\definecolor{hl}{RGB}{30,110,200}

% ---------------------------------------------------------------------
% Mini 3D solids for SOLUTIONS (items 54, 62, 66, 74): solid only, no
% washer, every picture 1.5 cm tall; a "big minus small" pair is 4.4 cm wide.
%  \solX[x unit cm]{R(x)}{a}{b}{R(a)}{R(b)}{x-axis end}{y-axis top}{curve colour}
%     solid swept out by y = R(x), a <= x <= b, about the x-axis
%  \solY[y unit cm]{R(y)}{c}{d}{R(c)}{R(d)}{x-axis end}{y-axis top}{curve colour}
%     solid swept out by x = R(y), c <= y <= d, about the y-axis
%     (R is written in terms of the letter x, which here stands for y)
% ---------------------------------------------------------------------
\newcommand{\solX}[9][0.5]{%
\begin{tikzpicture}[baseline=(current bounding box.center)]
\begin{axis}[mini3d, x={(#1cm,-0.06cm)}, y={(0.12cm,0.09cm)}, z={(0cm,0.4cm)}]
  \addplot3[cap3d, domain=0:1, y domain=0:360, samples=2, samples y=40] ({#3}, {x*(#5)*cos(y)}, {x*(#5)*sin(y)});
  \addplot3[solidsurf, opacity=0.5, point meta=z, domain=#3:#4, y domain=0:360, samples=24, samples y=40] ({x}, {(#2)*cos(y)}, {(#2)*sin(y)});
  \addplot3[cap3d, domain=0:1, y domain=0:360, samples=2, samples y=40] ({#4}, {x*(#6)*cos(y)}, {x*(#6)*sin(y)});
  \addplot3[#9, thick, domain=#3:#4, samples=40, samples y=1] ({x}, {0}, {#2});
  \draw[ax3] (axis cs:{min(0,#3)-0.3},0,0) -- (axis cs:#7,0,0);
  \draw[ax3] (axis cs:0,0,{-0.15*(#8)}) -- (axis cs:0,0,#8);
\end{axis}
\end{tikzpicture}}
\newcommand{\solY}[9][0.6]{%
\begin{tikzpicture}[baseline=(current bounding box.center)]
\begin{axis}[mini3d, x={(0.3cm,-0.03cm)}, y={(0.13cm,0.11cm)}, z={(0cm,#1cm)}]
  \addplot3[cap3d, domain=0:1, y domain=0:360, samples=2, samples y=40] ({x*(#5)*cos(y)}, {x*(#5)*sin(y)}, {#3});
  \addplot3[solidsurf, opacity=0.5, point meta=z, domain=#3:#4, y domain=0:360, samples=24, samples y=40] ({(#2)*cos(y)}, {(#2)*sin(y)}, {x});
  \addplot3[cap3d, domain=0:1, y domain=0:360, samples=2, samples y=40] ({x*(#6)*cos(y)}, {x*(#6)*sin(y)}, {#4});
  \addplot3[#9, thick, domain=#3:#4, samples=40, samples y=1] ({#2}, {0}, {x});
  \draw[ax3] (axis cs:0,0,{min(0,#3)-0.15*(#8)}) -- (axis cs:0,0,#8);
  \draw[ax3] (axis cs:-0.3,0,0) -- (axis cs:#7,0,0);
\end{axis}
\end{tikzpicture}}
\newcommand{\minisolo}[1]{\adjustbox{totalheight=1.5cm, max width=4.6cm}{#1}}
\newcommand{\minipair}[2]{\adjustbox{totalheight=1.5cm, max width=4.6cm}{#1\;\raisebox{0.15ex}{\Large$-$}\;#2}}
\newcommand{\minilabel}[1]{\raisebox{0.65cm}{\small #1}\;}


% Small fractions are allowed only in exponents and in the limits of an
% integral (which are superscripts/subscripts); everywhere else use \dfrac.
\newcommand{\dx}{\,\mathrm{d}x}
\newcommand{\dy}{\,\mathrm{d}y}
\newcommand{\evalat}[3]{\left[\,#1\,\right]_{#2}^{#3}}

% [item 78] right-aligned A-level style marks at the end of a question part
\newcommand{\mk}[1]{{\unskip\nobreak\hfil\penalty50\hskip1.5em\hbox{}\nobreak\hfill\textbf{[#1]}%
  \parfillskip=0pt\finalhyphendemerits=0\par}}
% [item 11] visually emphasised step labels in worked solutions
\newcommand{\step}[1]{\textcolor{teal!55!black}{\textbf{Step #1}}\quad}
% (a), (b), (c) part lists inside questions and solutions
\newenvironment{parts}{\begin{enumerate}[(a)]\setlength{\itemsep}{0.4ex}}{\end{enumerate}}
% SEAB-style solution layout (approved 10.5.8 template): working on the
% left, region sketch on the right (mini 3D pictures go under the sketch)
\newcommand{\solpair}[3][0.56]{\noindent\begin{minipage}[t]{#1\textwidth}\vspace{0pt}#2\end{minipage}\hfill
  \begin{minipage}[t]{\dimexpr0.98\textwidth-#1\textwidth\relax}\vspace{0pt}\centering #3\end{minipage}\par}
% minus-diagram style (drafts/minus_style_template.tex, items 29/31/62)
\pgfplotsset{minus/.style={axis lines=middle, width=5.0cm, height=4.2cm, ticks=none,
  xlabel={$x$}, ylabel={$y$}, every axis x label/.style={at={(ticklabel* cs:1)},anchor=west},
  every axis y label/.style={at={(ticklabel* cs:1)},anchor=south}, clip=false, axis on top,
  axis line style={gray!70}, every axis label/.append style={gray},
  every node/.append style={text=gray!80!black, font=\footnotesize}}}
\tikzset{hlreg/.style={fill=hl!35, draw=hl, thick}}

% =====================================================================
% prelim_fonts.tex: fonts of the real 9758 prelim papers (Word + MathType; pdffonts on
% ACJC 2026 P1/P2, ACJC 2025 P1, JPJC 2026, RVHS 2026, CJC/DHS/TJC/NYJC/HCI/EJC 2025):
%   text Times New Roman; cover/footer Arial; RI page footer Calibri (here Carlito);
%   maths = MathType: Times New Roman letters/digits (italic variables; upright digits,
%   function names, d, e, f), SymbolMT for the integral, sigma, pi and = + < >,
%   scripts 7 pt on 12 pt (about 58 %), tight operator spacing.
% Candidate C (teacher-approved 13 Oct 2026): the real Times New Roman and Arial
% (ttf-mscorefonts-installer) + URW Standard Symbols PS (fonts-urw-base35, a SymbolMT
% clone with the same Symbol encoding) for the integral, sigma, pi, = + < >;
% TeX Gyre Termes Math only for what MathType draws itself (radicals, stretchy brackets,
% fraction bars, arrows). House style: upright d (\mathrm{d}x), upright e and f.
% Loaded once by di_style / de_core (HW sheets) and prelim_style (Tests, solutions).
% =====================================================================
\ifdefined\prelimfontsloaded\endinput\fi
\def\prelimfontsloaded{1}
\edef\pfrestoreat{\catcode`\noexpand\@=\the\catcode`\@\relax}\makeatletter
% MathType script sizes: 7 pt superscripts on 12 pt text (58 %), 5 pt second level
\DeclareMathSizes{12}{12}{7}{5}
\DeclareMathSizes{11}{11}{6.5}{5}
\DeclareMathSizes{10.95}{10.95}{6.5}{5}
\DeclareMathSizes{10}{10}{6}{5}
\DeclareMathSizes{9}{9}{5.5}{5}
\DeclareMathSizes{14.4}{14.4}{8.5}{6}
\setmainfont{Times New Roman}
\setsansfont{Arial}
\newfontfamily\coverfont{Arial}
\newfontfamily\calibrifont{Carlito}
\usepackage{unicode-math}
\setmathfont{TeX Gyre Termes Math}
\setmathfont[range={up/{latin,Latin,num}}]{Times New Roman}
\setmathfont[range={it/{latin,Latin}}]{Times New Roman Italic}
\setmathfont[range={bfup/{latin,Latin,num}}]{Times New Roman Bold}
\setmathfont[range={"002B,"003D,"003C,"003E}]{Standard Symbols PS}
\newfontfamily\pf@sym{Standard Symbols PS}
\newfontfamily\pf@symi[FakeSlant=0.21]{Standard Symbols PS}
% MathType-like spacing (tighter than TeX around = + -)
\AtBeginDocument{\thickmuskip=2.5mu plus 1mu\medmuskip=1.5mu plus 1mu\thinmuskip=3mu
% Symbol integral and sigma (18 pt at 12 pt text, as MathType), slanted Symbol pi
 \renewcommand{\int}{\mathop{\vcenter{\hbox{\pf@sym\fontsize{\fpeval{1.5*\f@size}}{0pt}\selectfont\char"F2}}}\nolimits}%
 \renewcommand{\sum}{\mathop{\vcenter{\hbox{\pf@sym\fontsize{\fpeval{1.5*\f@size}}{0pt}\selectfont\char"E5}}}\limits}%
 \renewcommand{\pi}{\mathord{\text{\pf@symi\char"70}}}}
\pfrestoreat
% Times New Roman + MathType-style maths (as the prelim papers)
\usepackage{needspace}
\usepackage{zref-savepos,zref-abspage}\makeatletter\zref@addprop{savepos}{abspage}\makeatother
\usepackage{array}

% ---- shortlist toggle ------------------------------------------------
% \ifshortlist true  -> only the questions in SHORTLIST.md (+ challenges)
% A wrapper file sets it with \def\setshortlist{1} before \input-ing the
% main file; the main file can also be edited to say \shortlisttrue.
\newif\ifshortlist
% Task 5: the FINAL homework (= former shortlist) is the default build. HWn.tex -> HWn.pdf.
% A wrapper that sets \def\setfull{1} builds the archive with all original questions.
\ifdefined\setfull\else\shortlisttrue\fi

\makeatletter
% ---- REVIEW build: highlight amended content ---------------------------
% A wrapper sets \def\setreview{1}. In review mode amended text is shaded
% yellow and carries a hover note (PDF tooltip / sticky note); otherwise the
% macros print the content plainly.
%   \chg{note}{inline content}        short inline change (no line breaks), text mode
%   \chgm{note}{math}                 the same inside math mode
%   \begin{chgblock}{note} ... \end{chgblock}   paragraph/display-level change
\newif\ifreview
\ifdefined\setreview \reviewtrue \fi
% hover notes are raw PDF annotations (xdvipdfmx \special; no extra packages).
% Keep notes plain ASCII; balanced parentheses are fine.
\newsavebox{\chg@box}
\newcommand{\chg@ann}[4]{% width, height, depth, note -> invisible markup annotation with a hover note
  \special{pdf:ann width #1 height #2 depth #3 << /Type /Annot /Subtype /Square /Border [0 0 0] /BS << /W 0 >> /C [1 0.85 0] /CA 0.01
    /T (Review) /Subj (Amendment) /Contents (#4) >>}}
\newcommand{\chg@sticky}[1]{\special{pdf:ann width 12pt height 12pt depth 0pt << /Type /Annot /Subtype /Text /Name /Comment
    /C [1 0.85 0] /T (Review) /Subj (Amendment) /Contents (#1) >>}}
\newcommand{\chg}[2]{\ifreview\sbox\chg@box{\colorbox{yellow!60}{#2}}%
  \leavevmode\rlap{\chg@ann{\the\wd\chg@box}{\the\ht\chg@box}{\the\dp\chg@box}{#1}}\usebox\chg@box\else#2\fi}
% \chgm{note}{math}: the same, for use INSIDE math mode (display style kept)
\newcommand{\chgm}[2]{\ifreview\mathchoice{\chg{#1}{$\displaystyle#2$}}{\chg{#1}{$\textstyle#2$}}{\chg{#1}{$\scriptstyle#2$}}{\chg{#1}{$\scriptscriptstyle#2$}}\else#2\fi}
\newtcolorbox{chgbox}[1]{enhanced, breakable, colback=yellow!35, colframe=yellow!35, boxrule=0pt, sharp corners,
  left=2pt, right=14pt, top=1pt, bottom=1pt, before skip=2pt, after skip=2pt,
  overlay first and middle={}, overlay last={}, overlay unbroken and first={\node[anchor=north east, inner sep=1pt] at (frame.north east) {\chg@sticky{#1}\hbox to 12pt{\vrule width 0pt height 12pt}};}}
% \chgnote{note}: just the sticky-note icon (e.g. inside a figure); nothing outside review mode
\newcommand{\chgnote}[1]{\ifreview\chg@sticky{#1}\hbox to 12pt{\vrule width 0pt height 12pt}\fi}
\newenvironment{chgblock}[1]{\ifreview\begin{chgbox}{#1}\fi}{\ifreview\end{chgbox}\fi}

% ---- header ---------------------------------------------------------
\newcommand\di@title{}
\makeatletter
\newcommand{\ditle}[1]{\gdef\di@title{#1}}
\fancypagestyle{di}{\fancyhf{}%
  \lhead{\small \di@title}%
  \rhead{\small Compiled by Maximus Pung Jun}%
  \cfoot{\small\thepage}%
  \renewcommand{\headrulewidth}{0.4pt}}
\pagestyle{di}
% title block on page 1, as in his Typst sheets
\newcommand{\dititleblock}[1][]{\thispagestyle{di}%
  {\LARGE\scshape \di@title}\par\smallskip
  {Compiled by Maximus Pung Jun}\par\ifx&#1&\else\smallskip{\small\textcolor{black!60}{#1}}\par\fi\medskip}
\makeatother
\newcommand{\disection}[1]{\par\bigskip\noindent{\Large\scshape #1}\par\medskip}

% ---- questions --------------------------------------------------------
% \Qhead{label}{source tag}: question heading, tag in brackets on the right
\newcommand{\Qhead}[2]{\par\needspace{8\baselineskip}\bigskip\noindent
  {\large\bfseries #1}\hfill{\small\textcolor{black!60}{[#2]}}\par\smallskip}
% part lists with bold labels as he writes them: \pt{(a)}, \pt{(i)}, \pt{(b) (i)}
\newenvironment{qparts}[1][2.4em]{\begin{list}{}{\setlength\leftmargin{#1}%
  \setlength\labelwidth{\dimexpr#1-0.5em\relax}\setlength\labelsep{0.5em}%
  \setlength\itemsep{0.8ex}\setlength\topsep{0.6ex}\setlength\parsep{0.6ex}%
  \renewcommand\makelabel[1]{\textbf{##1}\hfil}}}{\end{list}}
\newcommand{\pt}[1]{\item[#1]}
% thin horizontal label for a challenge question
\newcommand{\challengebanner}{\par\bigskip\noindent\textcolor{black!50}{\rule[0.5ex]{\linewidth}{0.4pt}}\par
  \noindent{\large\scshape Challenge (optional)}\par}
% upright "units"
\newcommand{\units}[1]{\ \text{units}^{#1}}
% END line
\newcommand{\diend}[1][END OF WORKSHEET]{\par\vspace{2\baselineskip}\begin{center}\textbf{#1}\end{center}}

% ---- question registry (full / shortlist from one source) -------------
% \defQ{n}{tag}{body}   question n of the original sheet
% \defA{n}{answer body} its final answer, as printed in his answer section
% \putQ{n} / \putC{n}   print question n as the next (challenge) question
% \putA{n}              print its answer under the same new label
\newcounter{diq}
\makeatletter
\newcommand{\defQ}[3]{\expandafter\gdef\csname diQtag@#1\endcsname{#2}%
  \expandafter\long\expandafter\gdef\csname diQbody@#1\endcsname{#3}}
\newcommand{\defA}[2]{\expandafter\long\expandafter\gdef\csname diA@#1\endcsname{#2}}
\newcommand{\di@tag}[1]{\csname diQtag@#1\endcsname}% source tag only (no original Q number)
% A question never breaks across pages: it is measured and starts on a new page unless it
% fits whole on the rest of this one (only a question taller than a page may break, and then
% only between parts). Page positions of its start/end (zref labels q-s<n>/q-e<n> in the .aux)
% let the checker verify this.
\newsavebox{\di@qmeas}
\newcommand{\di@placeQ}[1]{\par\setbox\di@qmeas=\vbox{#1}%
  \ifdim\dimexpr\ht\di@qmeas+\dp\di@qmeas\relax<0.95\textheight
    \needspace{\dimexpr\ht\di@qmeas+\dp\di@qmeas+\bigskipamount\relax}\else\needspace{14\baselineskip}\fi
  \zsavepos{q-s\thediq}#1\par\zsavepos{q-e\thediq}}
\newcommand{\putQ}[1]{\stepcounter{diq}%
  \expandafter\xdef\csname diLab@#1\endcsname{\thediq}%
  \di@placeQ{\Qhead{\thediq}{\di@tag{#1}}\csname diQbody@#1\endcsname\par}}
\newcommand{\putC}[1]{\stepcounter{diq}%
  \expandafter\xdef\csname diLab@#1\endcsname{\thediq}%
  \expandafter\gdef\csname diChal@#1\endcsname{1}%
  \di@placeQ{\Qhead{\thediq\quad\normalsize\mdseries\scshape Challenge}{\di@tag{#1}}\csname diQbody@#1\endcsname\par}}
\newcommand{\putA}[1]{\par\needspace{4\baselineskip}%
  \begin{qparts}[2.4em]\item[{\csname diLab@#1\endcsname}]\csname diA@#1\endcsname\end{qparts}\medskip}
% ---- worked solutions: \defS{n}{body}, \putQS{n} prints question n + its solution ----
\newcommand{\defS}[2]{\expandafter\long\expandafter\gdef\csname diS@#1\endcsname{#2}}
\newcommand{\disolhead}{\par\medskip\noindent{\bfseries Solution}\par\nobreak\vspace{-0.6ex}\noindent\textcolor{black!35}{\rule{\linewidth}{0.4pt}}\par\nobreak}
\newcommand{\putQS}[1]{\putQ{#1}\disolhead\csname diS@#1\endcsname\par}
% source of a worked solution; visible only in REVIEW builds. #1: his / pub / pip
\newcommand{\solsrc}[2]{\ifreview\par\noindent\def\di@a{#1}\def\di@b{pip}%
  \ifx\di@a\di@b\colorbox{yellow!60}{\parbox{\dimexpr\linewidth-2\fboxsep-14pt}{\small\textbf{Pip-worked:} #2}}\chgnote{Pip-worked solution. No worked solution by the teacher and no published solution was found. Please check it.}%
  \else{\small\textcolor{black!55}{Source of this solution: #2}}\fi\par\smallskip\fi}
\makeatother
% compact sketch for worked solutions: \solfig[axis options]{plot commands}
\newcommand{\solfig}[2][]{\par\begin{center}\begin{tikzpicture}\begin{axis}[qfig, width=6.6cm, height=5cm, #1]#2\end{axis}\end{tikzpicture}\end{center}}
% a sketch added to an existing worked solution (teacher rule), shaded in REVIEW
\newenvironment{addedsketch}{\begin{chgblock}{Sketch added: teacher rule - every area or volume solution has a sketch.}}{\end{chgblock}}
% answer part list (compact)
\newenvironment{aparts}[1][2.4em]{\begin{qparts}[#1]\setlength\itemsep{0.3ex}}{\end{qparts}}

% ---- figure helpers ---------------------------------------------------------
\tikzset{rect/.style={fill=blue!10, draw=blue!55!black, thin}}
\pgfplotsset{qfig/.style={sketch, every node/.append style={font=\scriptsize}}}
% aligned working that starts on the same line as a part label: \pt{(iii)} \begin{topalign} ... \end{topalign}
\newenvironment{topalign}{\hfill$\displaystyle\begin{aligned}[t]}{\end{aligned}$\hfill\null\par\smallskip}
% one display that starts on the same line as a part label
\newcommand{\topdisplay}[1]{\hfill$\displaystyle #1$\hfill\null\par\smallskip}
